\subsection*{Kyle's rough notes}

\begin{itemize}

    \item Summary
    \begin{itemize}
        \item Develop approach for regressing non-Euclidean objects on high-dimensional predictors that does not impose parametric assumptions. 
        \item GFR (which generalizes linear regression) imposes strict assumptions. LFR (which generalizes nonparametric regression) suffers from the curse of dimensionality. DFR addresses these limitations -- i.e., handles high-dimensional predictors without compromising model flexibility -- by incorporating DNNs. 
        \item DFR Framework
        \begin{itemize}
            \item Learn $r$-dimensional representations of object-valued data using manifold learner ($\psi$; ISOMAP).
            \item Learn $r$ DNNs ($g_j$), each taking the predictors as input and outputting one coordinate of the manifold representation.
            \item Map conditional mean in representation space to object space via LFR ($v$), using the fitted values of the DNN as predictors. 
        \end{itemize}
    \end{itemize}

    \item Comments
    \begin{itemize}
    
        \item Comparison to SGM
        \begin{itemize}
            \item Analogies to SGM 
            \begin{itemize}
                \item ISOMAP like our encoder
                \item LFR like our decoder
                \item DNN like our latent statistical model
            \end{itemize}
            \item Their encoder (ISOMAP) aims to preserve geodesic distances in representation space while our encoder does not explicitly preserve distances. We instead encourage certain properties -- compactness, near-losslessness, Gaussianity -- and discover semantic structure.  So our models -- their DNN our statistical model -- are fit in spaces with fundamentally different geometries.  
        \end{itemize}

        \item Even though they compute conditional expectations in representation space then decode, their inferential targets are expectations in object space. This approximation, $\text{Decode}\{ \mathbb{E}(Z) \} \approx \mathbb{E} \{ \text{Decode}(Z) \}$, does not hold in general. In their setting, I think the approximation holds provided the representation space is geodesic-preserving. 
        \begin{itemize}
            \item Their approximation depends on ISOMAP reliably preserving geodesics. While they account for the bias introduced by ISOMAP in Theorem 2, the nature of this bias is not understood. 
            \item Choosing the latent dimension seems like an important question in their setting -- wrong choice of $r$ leads to representation that doesn't preserve geodesics, which leads to bad approximation. Do they describe how to choose $r$?
        \end{itemize}

        \item DNNs are learned independently in representation space. Not clear whether independent fitting is justified.

        \item DFR does not offer an inferential toolbox -- no testing, confidence intervals, etc. 

        \item LFR uses fitted values from DNNs to predict object-valued response. This decoding strategy is not cursed by dimensionality provided the latent dimension $r$ is small. However, LFR may struggle to stably decode for complex object data with moderate-to-large $r$.
        
    \end{itemize}

    \item Questions
    \begin{itemize}
        \item How are they choosing $r$? Are they just using $r = 2$?
        \item Is there something special about using LFR as the decoder?
    \end{itemize}

    \item Other Notes
    \begin{itemize}
        \item Faraway (2014) is another early predecessor to the encode-model-decode approach. The work ``scores'' (like ``embeds'') object-valued data by computing pairwise Euclidean distances, then applying multi-dimensional scaling. Worth citing in our work.
        \item Could be interesting to test manifold learners that explicitly preserve distances in our work. 
    \end{itemize}

    
\end{itemize}